



The training dataset for \MethodName{} consists of demonstration data for a wide range of tasks with many different robot platforms (both static and mobile, with single arm or bimanual) in diverse environments (in-house lab-like and home-like environments, and in-the-wild home environments), autonomous data from a large amount of policy evaluations, human interventions within policy rollouts, open-source robot datasets, egocentric human video data, and auxiliary non-robot data sources from the web, including object localization and attribute prediction, visual question answering, and text-only prediction. We also include video-language tasks including video captioning of in-house robot data and from the web.

In a significant departure from classic VLA training pipelines, we make heavy use of suboptimal robot data in training. This includes both lower quality demonstrations (failure episodes or success episodes with a substantial amount of mistakes) and data collected by prior versions of our models during model evaluation experiments\footnote{We exclude autonomous data collected in any generalization-focused evaluation task (including ones in Sec.~\ref{sec:experiments}) from training.}. For example, we use data collected by the \Piszs\ model during RL training as additional examples, effectively allowing \MethodName{} to distill their behavior. Incorporating the episode metadata into the context allows our model to effectively use all of this evaluation data and, as we will see in Sec.~\ref{sec:dexterity}, enables it to attain similar performance as models that are specialized for high performance on individual tasks with RL. This corresponds to a kind of ``distillation'' process, where the generalist \MethodName{} model can inherit the capabilities of RL-trained specialists. Suboptimal data also diversifies the possible states and scenarios in a given task and leads to stronger robustness, enabling the model to even sometimes outperform RL-trained or generally, single-task post-trained policies, in highly dexterous tasks.


